\documentclass[11pt,letterpaper]{article}
\usepackage[margin=1in]{geometry}
\usepackage{amsmath,amssymb,amsthm}
\usepackage{booktabs}
\usepackage{hyperref}
\usepackage{xcolor}
\usepackage{bm}

% ── Macros ─────────────────────────────────────────────────────────────────
\newcommand{\rcov}{r_{\mathrm{cov}}}
\newcommand{\nval}{n_{\mathrm{val}}}
\newcommand{\nmax}{n_{\mathrm{max}}}
\newcommand{\Zeff}{Z_{\mathrm{eff}}}
\newcommand{\kp}{k_p}
\newcommand{\cb}{c_b}
\newcommand{\asoft}{a_{\mathrm{soft}}}
\newcommand{\bchi}{b_\chi}
\newcommand{\Ang}{\text{\AA}}

\title{\textbf{A Generative Model for Atomic Static Dipole Polarizability}\\[4pt]
\large Descriptor Design, Fitting Methodology, and Validation (Z\,=\,1--120)}
\author{Formation Engine --- Atomistic Simulation Library}
\date{January 2025 \quad (Code: \texttt{atomistic/models/alpha\_model.hpp})}

\begin{document}
\maketitle

% ═══════════════════════════════════════════════════════════════════════════
\begin{abstract}
We present a 16-parameter generative model for the scalar isotropic static
dipole polarizability~$\alpha(Z)$ of neutral atoms, covering $Z=1$--$118$
with an experimental extension to $Z=120$.  The model requires no external
data files at runtime: all inputs are derived from the atomic number~$Z$
alone via closed-form descriptor functions (covalent radius, period, block,
electronegativity proxy, softness).  Parameters are fitted offline against
118 reference values from Miller~(1990) and Schwerdtfeger \&
Nagle~(2019) using iteratively-reweighted Huber regression with
group-monotonicity regularisation.  The overall RMS error is 19.1\% across
all 118 elements and $<$16\% for the commonly-simulated $Z=1$--$54$ subset.
The model feeds directly into the Thole-damped SCF induced-dipole solver
that computes polarization energies and forces in molecular dynamics.
\end{abstract}

% ═══════════════════════════════════════════════════════════════════════════
\section{Model Equation}
\label{sec:model}

The polarizability of element~$Z$ is predicted by
\begin{equation}
\boxed{\;
\alpha(Z) \;=\; \underbrace{\cb(Z)\;\kp(Z)\;\rcov(Z)^{3}\;
\frac{1 + \asoft\,s(Z)}{\chi(Z)^{\bchi}}}_{\alpha_{\text{base}}}
\;\cdot\; \underbrace{g_{\text{shell}}(Z)}_{\text{closure gate}}
\;\cdot\; \underbrace{g_{\text{rel}}(Z)}_{\text{relativ.\ stiffening}}
\;}
\label{eq:alpha}
\end{equation}
where the terms are:

\begin{center}
\begin{tabular}{lll}
\toprule
\textbf{Symbol} & \textbf{Meaning} & \textbf{Source} \\
\midrule
$\rcov(Z)$  & Covalent radius (\AA) & Pyykk\"o \& Atsumi (2009) \\
$s(Z)$      & Softness descriptor (dimensionless) & Eq.~\eqref{eq:soft} \\
$\chi(Z)$   & Electronegativity proxy & Eq.~\eqref{eq:chi} \\
$\kp(Z)$    & Period scaling factor ($p = 1 \ldots 7$) & Fitted \\
$\cb(Z)$    & Block correction ($b \in \{s,p,d,f\}$) & Fitted \\
$\asoft$    & Softness coupling coefficient & Fitted \\
$\bchi$     & Electronegativity suppression exponent & Fitted \\
$g_{\text{shell}}(Z)$ & Shell-closure suppression (Eq.~\ref{eq:gshell}) & Fitted \\
$g_{\text{rel}}(Z)$   & Relativistic stiffening (Eq.~\ref{eq:grel}) & Fitted \\
\bottomrule
\end{tabular}
\end{center}

\noindent
The model has \textbf{15 free parameters}: 7~period coefficients
$\kp^{(1)} \ldots \kp^{(7)}$, 4~block corrections
$\cb^{(s)},\cb^{(p)},\cb^{(d)},\cb^{(f)}$, 2~global exponents
$\asoft, \bchi$, and 2~gate amplitudes $c_{\text{shell}}, a_{\text{rel}}$.

\subsection{Physical Motivation}

The leading-order physics is the \emph{volume scaling law}: polarizability
grows as the cube of the electron-cloud radius.  For neutral atoms the
covalent radius~$\rcov$ is the best single-descriptor proxy for the
mean valence-shell extent.  The remaining terms correct for effects that
$\rcov^3$ alone cannot capture:

\begin{enumerate}
\item \textbf{Period scaling ($\kp$):}  Outer-shell electrons in higher
  periods are more diffuse per unit radius due to additional radial nodes.
  The fitted $\kp$ values decrease monotonically from period~1 to period~7,
  reflecting the diminishing ``volume efficiency'' as shielding increases.

\item \textbf{Electronegativity suppression ($\chi^{\bchi}$):}
  Electronegative atoms hold their electron cloud tighter, reducing
  deformability below what $\rcov^3$ predicts.  The denominator
  $\chi^{\bchi}$ implements this as a power-law suppression.

\item \textbf{Softness coupling ($1 + \asoft\,s$):}  The softness
  descriptor~$s$ measures how loosely-held the valence electrons are
  (high for alkali metals, low for noble gases).  The fitted
  $\asoft < 0$ \emph{reduces} $\alpha$ for soft atoms---a counter-intuitive
  result that compensates for the period scaling already being fit to
  alkali metals.

\item \textbf{Block correction ($\cb$):}  Transition metals ($d$-block)
  and lanthanides/actinides ($f$-block) have polarisable sub-valence
  electrons ($3d$, $4f$) whose spatial extent is not captured by $\rcov$.
  The block multiplier absorbs this orbital-type bias.
\end{enumerate}

% ═══════════════════════════════════════════════════════════════════════════
\section{Shell-Closure Gate}
\label{sec:shell}

\subsection{Why Noble Gases Blow Up (and How to Fix Without Hacks)}

Noble gases are closed-shell atoms: their electron clouds are stiff
and compact.  Polarizability is comparatively low and follows a
\emph{different curvature} than smooth volume scaling---shell closure
breaks the monotonic ``bigger radius $\Rightarrow$ more polarisable''
intuition that the base model assumes.

Before v2.6.2, the noble-gas RMS was $36.7\%$, with Ne at $+85\%$.
The existing $\chi$-based denominator in $\alpha_{\text{base}}$ partially
corrected for this (by assigning noble gases a high effective
electronegativity), but was insufficient because $\rcov^3$ still
dominates.

\subsection{The Shell-Closure Gate}

Noble gases should not be ``smoothed'' into the same curve as their
neighbours.  They must be \emph{gated}: a multiplicative penalty that
activates at or near shell closure.

We define a \emph{shell-closure proximity} descriptor:
\begin{equation}
\phi(Z) = \exp\!\left(-\frac{(18 - g(Z))^2}{2\sigma_{\text{shell}}^2}\right)
\label{eq:phi}
\end{equation}
where $g(Z)$ is the IUPAC group number ($g = 18$ for noble gases) and
$\sigma_{\text{shell}} = 1.8$.  This Gaussian is centred on group~18:
\begin{itemize}
  \item Noble gases ($g = 18$): $\phi = 1.0$ (full closure)
  \item Halogens ($g = 17$): $\phi \approx 0.86$ (near-closure)
  \item Chalcogens ($g = 16$): $\phi \approx 0.55$ (moderate)
  \item $g \leq 14$: $\phi < 0.15$ (negligible)
  \item $f$-block ($g = 0$): $\phi = 0$ by convention
\end{itemize}

The gate is then:
\begin{equation}
g_{\text{shell}}(Z) = 1 - c_{\text{shell}} \cdot \phi(Z)
\label{eq:gshell}
\end{equation}
where $c_{\text{shell}}$ is a single fitted amplitude.  The fitted
value $c_{\text{shell}} = 0.328$ means noble-gas polarizabilities are
suppressed by $\sim 33\%$ relative to the base model, while halogens
get a $\sim 28\%$ suppression---physically correct, as the halogen
anion $X^-$ (iso-electronic with the adjacent noble gas) has far
higher polarizability than neutral~$X$.

\paragraph{Design principle: gate, don't smooth.}
Smoothing residuals across a window that includes noble gases will
try to ``average out'' closure physics, creating a correction curve
that looks pretty and is physically wrong.  Noble gases are a
\emph{regime boundary}: the model must acknowledge the discontinuity
in electronic structure at shell closure, not paper over it.

% ═══════════════════════════════════════════════════════════════════════════
\section{Relativistic Stiffening}
\label{sec:rel}

\subsection{Why Mercury Spikes (and What Fixes It)}

Mercury ($Z = 80$) is the poster child for relativistic orbital
contraction: its $6s$ orbital contracts and stabilises, reducing
participation in polarizability and making Hg unusually ``inert''
compared to naive trend expectations.  This is the same physics
behind liquid mercury and its anomalous bonding.

Before v2.6.2, the model treated Hg like a normal heavy metal with
a fat, squishy outer cloud: predicted $7.79\;\text{\AA}^3$ vs.\ reference
$5.70\;\text{\AA}^3$ ($+37\%$).

\subsection{The Relativistic Stiffening Factor}

Relativistic contraction is \emph{not} monotonic with $Z$---it peaks
at specific electronic configurations.  We identify two hotspots:

\begin{enumerate}
\item \textbf{$6s$ inert-pair region} ($Z \approx 80$):
  Au($79$), Hg($80$), Tl($81$) all have anomalously compact $6s$
  orbitals due to strong nuclear shielding.

\item \textbf{$7p_{1/2}$ spin--orbit region} ($Z \approx 115$):
  The superheavy $p$-block (Nh--Og) experiences extreme $j$-dependent
  orbital splitting.  The $7p_{1/2}$ subshell contracts strongly,
  reducing the effective polarisable volume.
\end{enumerate}

We model these as a sum of two Gaussians:
\begin{equation}
f_{\text{rel}}(Z) = \exp\!\left(-\frac{(Z - 80)^2}{200}\right)
\;+\; \frac{1}{2}\,\exp\!\left(-\frac{(Z - 115)^2}{50}\right)
\label{eq:frel}
\end{equation}
and the gate is:
\begin{equation}
g_{\text{rel}}(Z) = 1 - a_{\text{rel}} \cdot f_{\text{rel}}(Z)
\label{eq:grel}
\end{equation}

The first Gaussian ($\sigma = 10$) yields $f_{\text{rel}} \approx 1.0$
at Hg, $\approx 0.7$ at Au/Tl, and decays to $< 0.1$ by Gd($64$)
and Po($84$).  The second Gaussian ($\sigma \approx 5$) peaks at
$0.5$ for Mc($115$) and covers Nh($113$)--Og($118$).

The fitted amplitude $a_{\text{rel}} = 0.054$ is physically small,
consistent with the expectation that relativistic effects are a
\emph{perturbative correction} to the volume-scaling base model, not
a dominant term.  The correction reduces Hg's error from $+37\%$ to
$+34\%$ and improves the overall period-6 RMS from $15.4\%$ to
$14.1\%$.

\paragraph{Why not a sigmoid?}
An earlier version used a sigmoid $\sigma(Z) = 1/(1 + e^{-(Z-76)/8})$.
This failed ($a_{\text{rel}} = 0$ in all fits) because it applies
uniform dampening to \emph{all} heavy atoms: Os, Ir, and Pt---which
fit well without correction---were degraded more than Hg was improved.
The double-Gaussian targets only the physical hotspots, allowing the
fitter to find a non-zero optimum.

\paragraph{Hg as a calibration anchor.}
If Hg's error improves, Au, Rn, and the $7p$ ``inert pair'' elements
(Nh--Og) typically improve as well, because the same $6s/7p_{1/2}$
contraction physics underlies all of them.  Hg should therefore be
monitored as a leading indicator in future model revisions.

% ═══════════════════════════════════════════════════════════════════════════
\section{Descriptor Definitions}
\label{sec:descriptors}

All descriptors are pure functions of $Z$.  No external database is
loaded at runtime.

\subsection{Covalent Radius $\rcov(Z)$}

Tabulated for $Z=1$--$118$ from Pyykk\"o \& Atsumi, \emph{Chem.\ Eur.\
J.}\ \textbf{15}, 186 (2009).  Stored as a compile-time constant array.
Units: \AA.

\subsection{Period and Block}

Standard IUPAC periodic table assignment:
\[
p(Z) \in \{1,2,3,4,5,6,7\}, \qquad
b(Z) \in \{s,\,p,\,d,\,f\}
\]
Lanthanides ($Z=57$--$71$) and actinides ($Z=89$--$103$) are assigned
$b = f$.  Group~0 (no IUPAC group number) is used for the $f$-block
insets; all other elements have standard group assignments $g = 1$--$18$.

\subsection{Active Valence Electron Count $\nval(Z)$}
\label{sec:nval}

\begin{equation}
\nval(Z) = \begin{cases}
  g           & g \in \{1,2\} \quad \text{($s$-block)} \\
  g - 10      & g \in \{13,\ldots,18\} \quad \text{($p$-block)} \\
  g           & g \in \{3,\ldots,12\} \quad \text{($d$-block)} \\
  3           & Z \in [57,71] \cup [89,103] \quad \text{($f$-block)}
\end{cases}
\label{eq:nval}
\end{equation}

\paragraph{$f$-block design decision.}
Lanthanide and actinide $4f/5f$ electrons are spatially contracted and
behave as core-like.  Setting $\nval = 3$ for all 30~elements prevents an
artificial monotonic gradient in $s(Z)$ across each series.  The actual
variation in polarizability within the lanthanide contraction is absorbed
by $\cb^{(f)}$ and $\rcov(Z)^3$.  This invariant is enforced by test~T1.

\subsection{Maximum Period Valence $\nmax(p)$}

\begin{equation}
\nmax(p) = \begin{cases}
  2  & p = 1 \\
  8  & p \in \{2,3\} \\
  18 & p \in \{4,5\} \\
  32 & p \in \{6,7\}
\end{cases}
\end{equation}

\subsection{Electronegativity Proxy $\chi(Z)$}
\label{sec:chi}

An Allred--Rochow-style estimate:
\begin{equation}
\chi(Z) = \frac{\nval^{\,\ast}(Z)}{\rcov(Z)^2} \times 0.359
\label{eq:chi}
\end{equation}
where the effective valence count $\nval^{\,\ast}$ equals $\nval(Z)$
for all elements except noble gases ($g=18$), which use
$\nval^{\,\ast} = \nmax(p)$ (full shell).  The normalisation constant
$0.359$ is chosen so that carbon ($\nval = 4$, $\rcov = 0.75\;\Ang$)
gives $\chi_{\mathrm{C}} \approx 2.55$ on the Pauling scale.

\paragraph{Noble gas electronegativity.}
Setting $\nval^{\,\ast} = \nmax$ for group~18 produces artificially high
$\chi$, which drives $\alpha$ down through the denominator of
Eq.~\eqref{eq:alpha}.  This is physically correct: noble gases have
compact, tightly-bound closed shells and small polarizabilities despite
moderate covalent radii.

\subsection{Softness $s(Z)$}
\label{sec:soft}

\begin{equation}
s(Z) = \frac{\nval(Z)\;/\;\nmax\bigl(p(Z)\bigr)}{\chi(Z)}
\label{eq:soft}
\end{equation}

Softness is high for diffuse, easily-deformable electron clouds (alkali
metals: $s_{\mathrm{K}} \approx 5.5$) and low for tight shells (fluorine:
$s_{\mathrm{F}} \approx 0.3$; neon: $s_{\mathrm{Ne}} \approx 0.03$).

% ═══════════════════════════════════════════════════════════════════════════
\section{Fitting Methodology}
\label{sec:fitting}

\subsection{Reference Data}

118 reference polarizabilities from two sources:
\begin{itemize}
\item \textbf{Miller (1990):} gas-phase experimental values for
  $Z=1$--$54$ (JACS~112, 8533).
\item \textbf{Schwerdtfeger \& Nagle (2019):} relativistic
  DHF/CCSD(T) values for $Z=21$--$118$ (Mol.\ Phys.~117, 1200).
\end{itemize}
Each entry carries a fitting weight $w_i \in [0.1, 2.0]$ encoding
reliability: $w = 2.0$ for high-priority simulation elements (H, C, N, O,
S), $w = 0.3$ for radioactive elements with theory-only values.

\subsection{Loss Function}

The objective is a weighted Huber loss in log-space plus a smoothness
penalty:
\begin{equation}
\mathcal{L}(\bm\theta) = \sum_{i=1}^{N} w_i \;\mathcal{H}_\delta\!\left(
  \ln\hat\alpha_i - \ln\alpha_i^{\mathrm{ref}}\right)
\;+\; \eta \sum_{(a,b)\in\mathcal{P}} \mathcal{H}_\delta\!\left(
  \Delta_{ab}^{\mathrm{pred}} - \Delta_{ab}^{\mathrm{ref}}\right)
\label{eq:loss}
\end{equation}
where $\mathcal{H}_\delta(x) = \tfrac{1}{2}x^2$ for $|x| \le \delta$
and $\delta(|x| - \tfrac{1}{2}\delta)$ otherwise, with $\delta = 0.3$.
The set $\mathcal{P}$ contains pairs of elements that are adjacent in the
same periodic-table group, sorted by period; $\Delta_{ab} = \ln\alpha_b -
\ln\alpha_a$ enforces monotonicity of trends down each group.

\paragraph{Why log-space.}  Polarizabilities span three orders of magnitude
($\alpha_{\mathrm{He}} = 0.205\;\Ang^3$ to
$\alpha_{\mathrm{Cs}} = 59.4\;\Ang^3$).  Log-space Huber loss treats
a~50\% error on helium as seriously as a~50\% error on caesium.

\subsection{Optimisation: Iterative Reweighted Coordinate Descent}

The 13~parameters are fitted by a two-loop scheme:

\begin{enumerate}
\item \textbf{Outer loop} (up to 30 rounds): after each inner
  convergence, update element weights $w_i$ based on residuals:
  \[
    w_i \leftarrow \min\!\left(w_{\max},\;
    w_i\bigl(1 + \alpha_{\mathrm{rw}} \cdot
    \max(0,\;|\epsilon_i|/\tau - 1)\bigr)\right)
  \]
  where $\epsilon_i = \ln\hat\alpha_i - \ln\alpha_i^{\mathrm{ref}}$,
  $\tau = 0.20$ is the tolerance target, $\alpha_{\mathrm{rw}} = 0.3$,
  and $w_{\max} = 10$.

\item \textbf{Inner loop} (up to 80 sweeps): coordinate descent over
  all 13~parameters.  Each parameter is optimised by a multi-pass grid
  search (4~passes $\times$ 100~grid points, geometrically refining the
  search window).
\end{enumerate}

\paragraph{Best-checkpoint selection.}
The weighted loss $\mathcal{L}$ monotonically increases across outer
iterations (because weights grow), but the \emph{unweighted} RMS error
may degrade after the optimal reweighting round.  The fitter tracks
the unweighted RMS at every outer iteration and returns the parameters
from the checkpoint with the lowest unweighted RMS, not the final
iterate.

% ═══════════════════════════════════════════════════════════════════════════
\section{Fitted Parameters (v2.6)}
\label{sec:params}

\begin{center}
\begin{tabular}{lrl}
\toprule
\textbf{Parameter} & \textbf{Value} & \textbf{Physical role} \\
\midrule
$\kp^{(1)}$ & 20.820 & Period 1 (H, He) \\
$\kp^{(2)}$ &  5.381 & Period 2 (Li--Ne) \\
$\kp^{(3)}$ &  5.117 & Period 3 (Na--Ar) \\
$\kp^{(4)}$ &  3.956 & Period 4 (K--Kr) \\
$\kp^{(5)}$ &  3.728 & Period 5 (Rb--Xe) \\
$\kp^{(6)}$ &  2.919 & Period 6 (Cs--Rn) \\
$\kp^{(7)}$ &  2.665 & Period 7 (Fr--Og) \\[4pt]
$\asoft$     & $-0.680$ & Softness coupling \\
$\bchi$      &  0.225 & Electronegativity exponent \\[4pt]
$\cb^{(s)}$  &  1.464 & $s$-block multiplier \\
$\cb^{(p)}$  &  0.973 & $p$-block multiplier \\
$\cb^{(d)}$  &  1.605 & $d$-block multiplier \\
$\cb^{(f)}$  &  1.745 & $f$-block multiplier \\[4pt]
$c_{\text{shell}}$ & 0.328 & Shell-closure suppression amplitude \\
$a_{\text{rel}}$   & 0.054 & Relativistic stiffening amplitude \\
\bottomrule
\end{tabular}
\end{center}

\noindent The monotonic decrease $\kp^{(1)} > \kp^{(2)} > \cdots >
\kp^{(7)}$ reflects increasing nuclear shielding: each additional shell
of core electrons makes the volume-to-polarizability conversion less
efficient.  The $d$-block multiplier $\cb^{(d)} \approx 1.6$ confirms
that partially-filled $d$ shells contribute significant
polarisable density beyond what $\rcov^3$ alone predicts.
The shell-closure amplitude $c_{\text{shell}} \approx 0.33$ suppresses
noble-gas predictions by a third, while the relativistic amplitude
$a_{\text{rel}} \approx 0.05$ provides a targeted $\sim 5\%$ correction
at the Hg hotspot.

% ═══════════════════════════════════════════════════================================================================
\section{Validation Results}
\label{sec:validation}

\subsection{Global Accuracy}

\begin{center}
\begin{tabular}{lrr}
\toprule
\textbf{Metric} & \textbf{Z = 1--118} & \textbf{Z = 1--54} \\
\midrule
RMS error              & 17.5\% & 14.8\% \\
Max error              & 64.5\% (Lv) & 57.4\% (Ne) \\
Monotonicity breaks    & 10/70 pairs & 5/35 pairs \\
Elements $>$20\% error & 23/118 & 8/54 \\
\bottomrule
\end{tabular}
\end{center}

\subsection{Per-Block RMS}

\begin{center}
\begin{tabular}{lrr}
\toprule
\textbf{Block} & \textbf{RMS} & \textbf{Physical note} \\
\midrule
$s$-block & 19.0\% & Alkali/alkaline-earth: large $\alpha$, well constrained \\
$p$-block & 19.9\% & Noble-gas suppression now active ($g_{\text{shell}}$) \\
$d$-block & 15.5\% & Transition metals: $\cb^{(d)}$ absorbs $d$-shell effect \\
$f$-block & 16.0\% & Lanthanides/actinides: fixed $\nval=3$ works well \\
\bottomrule
\end{tabular}
\end{center}

\subsection{Improvement from v2.6.0 to v2.6.2}

\begin{center}
\begin{tabular}{lrrr}
\toprule
\textbf{Regime} & \textbf{v2.6.0} & \textbf{v2.6.2} & \textbf{Improvement} \\
\midrule
Overall RMS           & 19.1\% & 17.5\% & $-8.4\%$ \\
Noble gases RMS       & 36.7\% & 27.3\% & $-25.6\%$ \\
Halogens RMS          & 24.1\% & 14.2\% & $-41.1\%$ \\
Period 6 RMS          & 15.4\% & 14.1\% & $-8.4\%$ \\
Hg error              & $+36.6\%$ & $+34.0\%$ & $-7.1\%$ \\
Ne error              & $+85.0\%$ & $+57.4\%$ & $-32.5\%$ \\
Max error             & 85.0\% & 64.5\% & $-24.1\%$ \\
\bottomrule
\end{tabular}
\end{center}

\subsection{Known Failure Modes}

\begin{enumerate}
\item \textbf{Neon} ($\alpha_{\mathrm{ref}} = 0.396\;\Ang^3$, error
  $+85\%$): Ne has an anomalously small polarizability relative to its
  covalent radius.  The model's $\chi$-based suppression is insufficient
  because $\rcov(\mathrm{Ne}) = 0.67\;\Ang$ overestimates the true
  electron-cloud extent for this closed-shell atom.

\item \textbf{Nobelium} ($\alpha_{\mathrm{ref}} = 17.5\;\Ang^3$, error
  $+48\%$): No has strong relativistic contraction of the $7s$ orbital
  that reduces $\alpha$ below the $\rcov^3$ expectation.  Our
  non-relativistic descriptors cannot capture this.

\item \textbf{Superheavy $p$-block} (Lv, Ts: errors $>$40\%):
  Extreme spin--orbit coupling in the $7p_{1/2}$ / $7p_{3/2}$ manifold
  changes the effective orbital extent.  Relativistic CCSD(T) reference
  values carry significant uncertainty themselves.
\end{enumerate}

% ═══════════════════════════════════════════════════════════════════════════
\section{Integration with the SCF Polarization Solver}
\label{sec:scf}

The alpha model provides the atomic polarizabilities
$\alpha_i = \alpha(Z_i)$ that enter the Applequist--Thole SCF equations:
\begin{align}
\bm\mu_i &= \alpha_i \, \mathbf{E}_i^{\mathrm{loc}} \label{eq:scf1} \\
\mathbf{E}_i^{\mathrm{loc}} &= \mathbf{E}_i^{\mathrm{perm}}
  + \sum_{\substack{j \ne i \\ r_{ij} > r_{\mathrm{excl}}}}
  f(u_{ij})\,\mathbf{T}_{ij}\,\bm\mu_j \label{eq:scf2}
\end{align}
where $\mathbf{T}_{ij}$ is the dipole--dipole interaction tensor and
$f(u)$ is the Thole exponential damping function:
\begin{equation}
f(u) = 1 - \left(1 + \tfrac{a\,u}{2}\right)e^{-au}, \qquad
u_{ij} = \frac{r_{ij}}{(\alpha_i\,\alpha_j)^{1/6}}, \qquad
a = 2.6
\label{eq:thole}
\end{equation}

\paragraph{Unit convention.}
All electric fields ($\mathbf{E}^{\mathrm{perm}}$ from charges,
$\mathbf{T}_{ij}\bm\mu_j$ from dipoles) are computed in Gaussian units
($e/\Ang^2$) without the Coulomb constant $k_e$.  The constant
$k_e = 332.0636\;\mathrm{kcal\cdot\Ang/(mol\cdot e^2)}$ appears only in
the energy formula:
\begin{equation}
U_{\mathrm{pol}} = -\frac{1}{2}\,k_e \sum_{i=1}^{N}
  \bm\mu_i \cdot \mathbf{E}_i^{\mathrm{perm}}
\label{eq:upol}
\end{equation}

\paragraph{Exclusion radius.}
Pairs with $r_{ij} < r_{\mathrm{excl}} = 1.6\;\Ang$ are excluded from
both the permanent-charge driving field and the dipole--dipole tensor.
This is equivalent to standard 1-2/1-3 exclusions in AMOEBA without
requiring bond topology.

% ═══════════════════════════════════════════════════════════════════════════
\section{Experimental Extension: Z\,=\,119--120}
\label{sec:ext}

Elements 119 (ununennium, Uue) and 120 (unbinilium, Ubn) have not been
synthesised in measurable quantities.  An optional extension module
(\texttt{alpha\_model\_ext.hpp}) provides:

\begin{itemize}
\item \textbf{Literature values}: $\alpha_{119} = 165\;\Ang^3$,
  $\alpha_{120} = 109\;\Ang^3$ from 4-component DHF calculations
  (Schwerdtfeger \& Nagle, 2019).
\item \textbf{Generative prediction}: the standard model equation using
  extrapolated period-8 descriptors ($\rcov(119) = 2.60\;\Ang$,
  $\rcov(120) = 2.21\;\Ang$ from Pyykk\"o 2015).  The period-7
  coefficient $\kp^{(7)}$ is reused as a floor.
\end{itemize}

\noindent The generative predictions ($\alpha_{119}^{\mathrm{model}}
\approx 103\;\Ang^3$, $\alpha_{120}^{\mathrm{model}} \approx 52\;\Ang^3$)
underestimate the literature values, confirming that relativistic
$8s$-contraction effects would require a period-8-specific $\kp^{(8)}$
coefficient that cannot be fitted without training data.

% ═══════════════════════════════════════════════════════════════════════════
\section{Software Architecture}
\label{sec:arch}

\begin{center}
\begin{tabular}{ll}
\toprule
\textbf{File} & \textbf{Role} \\
\midrule
\texttt{atomic\_descriptors.hpp}     & Pure-function descriptor layer ($\rcov$, $\chi$, $s$, period, block) \\
\texttt{alpha\_model.hpp}            & Model equation + baked parameters (13 coefficients) \\
\texttt{alpha\_model\_ext.hpp}       & Optional Z=119--120 extension \\
\texttt{polarization\_scf.cpp}       & Thole-damped SCF solver (consumes $\alpha_i$) \\
\texttt{fit\_alpha\_model.cpp}       & Offline fitter (reads CSV, writes JSON + updates baked params) \\
\texttt{polarizability\_ref.csv}     & Reference data (118 entries, not shipped to production) \\
\texttt{test\_alpha\_model.cpp}      & Descriptor invariants (T1--T4, 56 checks) \\
\texttt{test\_nuclear\_chemical.cpp} & Phase 6 validation (F1--F5, 24 checks) \\
\bottomrule
\end{tabular}
\end{center}

\noindent Key design constraint: the descriptor layer and model equation
are \emph{read-only at runtime}.  All tuning happens offline in the fitter;
the production binary contains only the baked 13-parameter struct.

% ═══════════════════════════════════════════════════════════════════════════
\section*{References}

\begin{enumerate}
\item[{[1]}] K.~J.~Miller, ``Additivity methods in molecular
  polarizability,'' \emph{J.\ Am.\ Chem.\ Soc.}\ \textbf{112}, 8533
  (1990).

\item[{[2]}] P.~Schwerdtfeger and J.~K.~Nagle, ``2018 Table of static
  dipole polarizabilities of the neutral elements in the periodic table,''
  \emph{Mol.\ Phys.}\ \textbf{117}, 1200 (2019).

\item[{[3]}] P.~Pyykk\"o and M.~Atsumi, ``Molecular single-bond covalent
  radii for elements 1--118,'' \emph{Chem.\ Eur.\ J.}\ \textbf{15}, 186
  (2009).

\item[{[4]}] J.~Applequist, J.~R.~Carl, and K.-K.~Fung, ``An atom dipole
  interaction model for molecular polarizability,'' \emph{J.\ Am.\ Chem.\
  Soc.}\ \textbf{94}, 2952 (1972).

\item[{[5]}] B.~T.~Thole, ``Molecular polarizabilities calculated with a
  modified dipole interaction,'' \emph{Chem.\ Phys.}\ \textbf{59}, 341
  (1981).

\item[{[6]}] P.~Pyykk\"o, ``Additive covalent radii for single-, double-,
  and triple-bonded molecules and tetrahedrally bonded crystals: a
  summary,'' \emph{J.\ Phys.\ Chem.\ A}\ \textbf{119}, 2326 (2015).

\item[{[7]}] J.~W.~Ponder and D.~A.~Case, ``Force fields for protein
  simulations,'' \emph{Adv.\ Protein Chem.}\ \textbf{66}, 27 (2003).
\end{enumerate}

% ═══════════════════════════════════════════════════════════════════════════
\section{Nuclear Stability Integration}
\label{sec:nuclear}

The polarizability model is tied into a nuclear stability module
(\texttt{nuclear\_stability.hpp}) that silently annotates each prediction
with a confidence score derived entirely from atomic structure.

\subsection{Ab-Initio Stability from $Z$}

Every call to \texttt{alpha\_predict\_checked(Z)} runs a silent stability
check using three physics engines—no external decay data:

\begin{enumerate}
\item \textbf{Semi-empirical mass formula} (Bethe--Weizs\"acker, 5
  parameters + shell correction):  Predicts binding energy $B(Z,A)$,
  most stable mass number $A^*$, and the fissility parameter $Z^2/A$.

\item \textbf{Nuclear shell model} (magic numbers 2, 8, 20, 28, 50, 82,
  126):  Gaussian shell correction to $B$ enhances stability at closed
  proton/neutron shells.

\item \textbf{Geiger--Nuttall law} (Viola--Seaborg parameterisation):
  Estimates $\alpha$-decay half-life from the $Q$-value
  $Q_\alpha = B(Z{-}2,A{-}4) + B(2,4) - B(Z,A)$.
\end{enumerate}

\subsection{Stability Classes and Confidence}

\begin{center}
\begin{tabular}{llrl}
\toprule
\textbf{Class} & \textbf{Criterion} & \textbf{$\alpha$-conf.} & \textbf{Example} \\
\midrule
Stable          & $\exists$ stable isotope & 1.00  & C, Fe, Pb \\
PrimordialLong  & $t_{1/2} > 4.5$\,Gy      & 0.85  & Th, U, Bi \\
Radioactive     & $t_{1/2} > 1$\,s          & 0.47--0.90 & Tc, Pm, Cf \\
VeryShortLived  & $t_{1/2} \in [1\,\mu$s, 1\,s) & 0.10--0.30 & --- \\
Superheavy      & $t_{1/2} < 1$\,ms or $Z \geq 104$ & 0.02--0.25 & Cn, Lv, Og \\
\bottomrule
\end{tabular}
\end{center}

\noindent
The confidence score gates the AlphaCalibrator's acceptance threshold:
stable elements use the strict $1\%$ gate; superheavy elements widen to
$20\%$ (acknowledging that their polarizabilities are purely theoretical).

\subsection{Known Limits}

\begin{itemize}
\item \textbf{He-4}: Doubly magic ($Z{=}2$, $N{=}2$).  Hardcoded $A^* = 4$
  because the SEMF liquid-drop model cannot resolve this light nucleus.

\item \textbf{Tc, Pm}: Correctly identified as having no stable isotopes
  despite $Z < 83$, classified as Radioactive ($\beta^-$ decay).

\item \textbf{Z=104--118}: Classified Superheavy with confidence
  $\leq 0.25$.  Half-lives are clamped at $\leq 10^5$\,s to prevent
  SEMF over-prediction of stability for transactinides.
\end{itemize}

\end{document}
